\chapter{知识的整合}
\label{chap:ket-integration}

花卉识别系统已经保存了两个分别训练的分类器：一个主要接触清晰的近景照片，另一个接触更多低清晰度图像。它们识别相同的物种，却不总在相同样本上判断正确。系统还保存了若干识别颜色、形态和病斑的适配模块。现在的问题不只是选出一个最好的来源，而是怎样让这些已经获得的能力共同工作：可以同时询问两个分类器，可以把它们的参数合成一份，也可以训练另一个模型承接它们的行为。这些操作究竟保留了什么，又在哪些条件下可能失效？

本章以\term{知识整合}（knowledge integration）指将分别获得的模型、模块或其他知识载体共同用于目标行为的过程。这里的知识是支持计算和学习的信息，不意味着每一组参数都对应一条可分离的语义事实。为比较三种机制，先固定花卉分类任务、类别定义和独立评价样本，再改变来源能力进入结果的方式；需要研究新的属性或指令组合时，会明确改变目标任务。全文数值例子均为教学构造，不是实验测量。

图~\ref{fig:ket-integration-routes}把三种关系并列画出。\term{模型组合}（model composition）保留来源计算单元，在预测时共同调用；\term{参数融合}（model merging）读取来源权重，构造一份新的参数；\term{知识蒸馏}（knowledge distillation）让来源提供监督，通过训练使承接模型学会相应行为。三者不是必须依次执行的阶段：一个组合系统也可以成为蒸馏教师。区别在于来源以计算、参数还是训练信号的形式进入结果。

\begin{table*}[htbp]
  \centering
  \small
  \begin{tabularx}{\linewidth}{@{}p{22mm}XXXX@{}}
    \toprule
    路线 & 来源访问 & 额外材料与训练 & 部署形态 & 首要失败边界 \\
    \midrule
    模型组合 & 运行时可执行来源；模块组合还需中间接口 & 聚合器或路由器可需验证数据和训练 & 多来源、路由或共享模块继续存在 & 接口错位、共同错误、路由失配与重复计算 \\
    参数融合 & 可读取且可对应的来源参数与状态 & 可不用数据，也可用输入、标签或激活确定对应与系数 & 通常一份参数，可保留多个头 & 坐标不对应、跨越高损失区与来源能力冲突 \\
    知识蒸馏 & 可查询输出；表示蒸馏还需内部激活 & 覆盖目标范围的输入、学生训练和模型选择 & 只部署学生，教师可离线移除 & 查询覆盖不足、教师偏差与学生容量限制 \\
    \bottomrule
  \end{tabularx}
  \caption{三条整合路线的统一决策坐标。路线可以串联，但来源访问、训练成本、部署状态和失败边界须分别记账；“最终只有一个模型”不表示整合过程无数据或无成本。}
  \label{tab:ket-integration-route-choice}
\end{table*}

\begin{figure*}[htbp]
  \centering
  % Three independent integration routes; schematic, no empirical data.
% Final canvas: 164 mm x 131 mm. Double outlines mark deployed units.
% Solid: forward computation; dashed: training; dotted: parameter reading.
\begin{tikzpicture}[
  x=1mm,y=1mm,font=\figurefont\footnotesize,text=BookInk,
  box/.style={rectangle,draw=FigureModel,fill=FigureModel!7,
    line width=0.8pt,minimum width=25mm,minimum height=9mm,
    inner sep=1mm,align=center},
  deployed/.style={box,double,double distance=0.7pt},
  input/.style={box,draw=FigureInput,fill=FigureInput!7,
    minimum width=15mm},
  result/.style={box,draw=FigureOutput,fill=FigureOutput!7,
    minimum width=15mm},
  loss/.style={box,draw=FigureLoss,fill=FigureLoss!7},
  flow/.style={-{Latex[length=2mm,width=1.3mm]},draw=BookInk,
    line width=0.8pt,rounded corners=1mm},
  train/.style={flow,dashed,draw=FigureLoss},
  weights/.style={flow,dotted,draw=BookBlue},
  construct/.style={flow,draw=BookBlue},
  note/.style={font=\figurefont\footnotesize,text=BookMuted,inner sep=1mm}
]
  \path[use as bounding box] (0,-13) rectangle (164,118);
  \node[anchor=west,font=\figurefont\small] at (1,113) {(a) 保留来源计算};
  \node[input] (xa) at (12,94) {$\bm{x}$};
  \node[deployed] (a1) at (47,101) {来源$f_1$};
  \node[deployed] (a2) at (47,87) {来源$f_2$};
  \node[deployed] (agg) at (99,94) {聚合器};
  \node[result] (ya) at (150,94) {预测};
  \draw[flow] (xa.east) -- (25,94) |- (a1.west);
  \draw[flow] (25,94) |- (a2.west);
  \draw[flow] (a1.east) -- (76,101) |- (agg.west);
  \draw[flow] (a2.east) -- (76,87) |- (agg.west);
  \draw[flow] (agg.east) -- (ya.west);
  \draw[BookRule,line width=0.6pt] (0,78)--(164,78);

  \node[anchor=west,font=\figurefont\small] at (1,73) {(b) 读取来源参数};
  \node[box] (b1) at (34,60) {来源$f_1$\\$\bm{\theta}_1$};
  \node[box] (b2) at (34,45) {来源$f_2$\\$\bm{\theta}_2$};
  \node[box] (rule) at (76,53) {融合规则};
  \node[deployed] (merged) at (122,53) {融合模型};
  \node[result] (yb) at (151,53) {预测};
  \node[input] (xb) at (92,37) {$\bm{x}$};
  \draw[weights] (b1.east) -- (53,60) |- (rule.west);
  \draw[weights] (b2.east) -- (53,45) |- (rule.west);
  \draw[construct] (rule.east) -- node[note,above] {构造新参数} (merged.west);
  \draw[flow] (xb.east) -| (merged.south);
  \draw[flow] (merged.east) -- (yb.west);
  \draw[BookRule,line width=0.6pt] (0,29)--(164,29);

  \node[anchor=west,font=\figurefont\small] at (1,24) {(c) 用来源监督训练};
  \node[input] (xc) at (12,8) {$\bm{x}$};
  \node[box] (c1) at (45,15) {来源$f_1$};
  \node[box] (c2) at (45,1) {来源$f_2$};
  \node[loss] (kd) at (85,8) {蒸馏约束};
  \node[deployed] (student) at (124,8) {学生};
  \node[result] (yc) at (152,8) {预测};
  \draw[flow] (xc.east) -- (24,8) |- (c1.west);
  \draw[flow] (24,8) |- (c2.west);
  \draw[train] (c1.east) -- (65,15) |- (kd.west);
  \draw[train] (c2.east) -- (65,1) |- (kd.west);
  \draw[train] (kd.east) -- node[note,above] {更新} (student.west);
  \draw[flow] (student.east) -- (yc.west);
  \draw[flow] (student.north) -- (124,21) -| (kd.north);
  \node[note] at (105,23) {学生输出};
  \draw[flow] (xc.south) -- (12,-9) -| (student.south);
  \node[note,fill=white] at (86,-9) {同一输入};
\end{tikzpicture}

  \caption{同一对来源模型的三种整合路径。黑色实线为前向计算，虚线为训练约束或参数更新，点线只表示从来源读取参数，蓝色实线表示融合规则构造新参数；带双重轮廓的单元在部署时仍需保存并运行。组合保留两个来源与聚合器，融合只执行构造出的模型，蒸馏只执行训练后的学生。各行彼此独立，未表达三种操作的先后顺序。}
  \label{fig:ket-integration-routes}
\end{figure*}

第~\ref{chap:ket-reuse-selection}~章提供来源可用性与复用收益的判断，第~\ref{chap:ket-update-maintenance}~章讨论单个系统的更新与维护。本章从已经形成的来源出发，考察它们怎样进入组合结果；第~\ref{sec:ket-sharing-computation}~节则考察多个任务怎样共同更新共享计算。是否还需训练组合器或学生不是两者的分界，来源能力是否已经给定、哪些任务共同改变共享量才是关键。训练材料、标签和记录本身怎样产生与校验，仍属于第~\ref{part:efficient-knowledge-use}~部分的知识获取主题。

\section{基于模型组合的知识整合}
\label{sec:ket-integration-composition}

保留来源能够避免立即改写它们的内部参数，但还必须确定组合接口。若各来源都回答“这是什么花”，整合的是同义预测；若一个来源识别花瓣形态，另一个依据形态判断物种，整合的就是中间计算。完整模型和小型适配器都可能承担其中一种角色，计算单元的大小不是分类标准。

在每种接口内部，还要先约定方案服务的任务和输入范围。在该范围内始终使用同一组来源、同一系数或同一连接，称为固定组合；随任务或输入改变选择与组合，称为条件组合。固定权重可以从数据中学得，条件规则也可以人工给定。是否经过训练、是否在部署前确定，均不能代替固定与条件的区分。

\subsection{预测结果的集成}
\label{sec:ket-integration-prediction}

把多个答案合在一起之前，需要确认它们在回答同一个问题。设第$k$个分类器输出$p_k(y\mid\bm{x})$，类别集合为$Y=\{1,\ldots,c\}$。两个长度相同的向量，可能分别按“月季、百合、菊花”和“百合、菊花、月季”排序；不先重排类别，平均得到的数值便没有所声明的概率含义。多标签概率与互斥类别概率、物种概率与病害概率也不能直接对应。即使类别相同，来源对同一数值置信度的实际正确率仍可能不同，需要在与目标使用范围相符的材料上检查校准。

第~\ref{chap:ensemble-learning}~章已经讨论集成模型怎样训练成员。本节把成员视为已选定的来源，关心它们的答案怎样共同使用、运行时调用哪些模型，以及这种操作能支持何种结论。

\subsubsection{基于固定规则的预测集成}
\label{sec:ket-integration-fixed-prediction}

多数投票只保留每个来源的最终类别，再选择票数最多的类别；两个来源可能一票对一票，必须另定平局规则。概率平均保留每个来源对全部类别的相对倾向。更一般的固定权重规则为
\begin{equation}
  \bar p(y\mid\bm{x})=\sum_{k=1}^{K}\alpha_k p_k(y\mid\bm{x}),
  \qquad \alpha_k\geq 0,\quad \sum_{k=1}^{K}\alpha_k=1.
  \label{eq:ket-integration-fixed-average}
\end{equation}
这里$K$为来源数，系数在约定范围内不随输入改变。非负且和为一保证$\bar p$仍是概率分布，却不保证它校准良好或判断正确。可以预先在独立验证集上按损失选择$\bm{\alpha}$，随后冻结系数；用于选择的验证成绩不能替代最终测试成绩。

\begin{example}[互补错误与共同错误]
  \label{ex:ket-integration-weighted-prediction}
  为便于手算，把花卉任务简化为二分类，$y=1$表示月季，阈值取$0.5$，等于阈值时判为正类。在同一组四张照片上，真实标签及两个来源对正类的概率为
  \[
    (y_1,y_2,y_3,y_4)=(1,0,1,0),\qquad
    \begin{aligned}
      (p_1(\bm{x}_i))_{i=1}^{4}&=(0.9,0.4,0.6,0.6),\\
      (p_2(\bm{x}_i))_{i=1}^{4}&=(0.4,0.1,0.9,0.3).
    \end{aligned}
  \]
  两个模型各错一张，正确率均为$3/4$。取$\alpha_1=0.6,\alpha_2=0.4$，得到
  \[
    (\bar p(\bm{x}_i))_{i=1}^{4}
    =(0.70,0.28,0.72,0.48).
  \]
  例如第四张的计算是$0.6\times0.6+0.4\times0.3=0.48$，组合后四张全部判断正确。若保持第一个模型、标签和权重不变，只把第二个模型的概率改为$(0.8,0.2,0.7,0.7)$，第二个模型仍有$3/4$的正确率，但两个来源都错第四张。此时组合概率变为$(0.86,0.32,0.64,0.64)$，错误没有消失。相同的成员正确率和数量，并不决定相同的集成收益。
\end{example}

错误之间的关系可以在平方损失下精确展开。考虑实值回归$y=f_\star(\bm{x})+\varepsilon$，其中$f_\star(\bm{x})=\mathbb{E}[y\mid\bm{x}]$，$\mathbb{E}[\varepsilon\mid\bm{x}]=0$。固定一个输入$\bm{x}$，把训练样本和训练随机性记为随机来源，假设二阶矩存在，测试噪声与训练过程条件独立。令$e_k=\hat f_k(\bm{x})-f_\star(\bm{x})$，$b_k=\mathbb{E}[e_k]$，$\Sigma_{k,j}=\operatorname{Cov}(e_k,e_j)$。由于权重固定且和为一，平均预测相对回归函数的误差为$\bar e=\sum_k\alpha_ke_k$，于是
\begin{equation}
  \begin{aligned}
    \mathbb{E}[\bar e^2]
    &=\biggl(\sum_k\alpha_k b_k\biggr)^2+
      \mathbb{E}\biggl[\biggl(\sum_k\alpha_k(e_k-b_k)\biggr)^2\biggr]\\
    &=\biggl(\sum_k\alpha_k b_k\biggr)^2+
      \sum_k\alpha_k^2\Sigma_{k,k}
      +2\sum_{k<j}\alpha_k\alpha_j\Sigma_{k,j}.
  \end{aligned}
  \label{eq:ket-integration-bias-covariance}
\end{equation}
第一步把平方误差分成均值平方与中心化误差的方差；第二步展开平方，每对不同来源产生一个协方差项。若比较对带噪声标签$y$的平方损失，还须加上$\operatorname{Var}(\varepsilon\mid\bm{x})$。输入不固定时，再对目标输入分布取期望。

若每个来源无偏、方差均为$\sigma^2$，两两相关系数均为$\rho$，等权平均便有
\begin{equation}
  \operatorname{Var}(\bar e)
  =\frac{K\sigma^2+K(K-1)\rho\sigma^2}{K^2}
  =\sigma^2\left(\rho+\frac{1-\rho}{K}\right).
  \label{eq:ket-integration-equal-correlation}
\end{equation}
当$\rho=1$时增加来源不降方差；当$\rho=0$时方差变为$\sigma^2/K$。这里的$\rho$必须使协方差矩阵有效，$K>1$时等相关结构要求$\rho\geq-1/(K-1)$。共同偏差也不会因为平均自动消失。这是固定权重、平方损失和所列随机性条件下的关系，不是分类准确率必然提高的证明；分类还取决于概率跨越决策边界的方式。

平均未归一化分数与平均概率也不同。设来源分数为$\bm{a}_k$，概率为$\operatorname{softmax}(\bm{a}_k)$，一般有$\operatorname{softmax}(\sum_k\alpha_k\bm{a}_k)\neq\sum_k\alpha_k\operatorname{softmax}(\bm{a}_k)$。前者相当于对严格正的类别概率作加权几何平均后重新归一化，后者是算术平均；投票则丢弃了概率幅度。分数尺度偏大的来源可能支配分数平均，而高置信度也可能只是过度自信，不能据此直接提高权重。

固定集成通常要运行全部$K$个来源。独立模型的参数存储近似相加，总计算量也近似相加；并行执行可能缩短墙钟延迟，却受设备数量、通信和最慢来源制约，串行执行则累积延迟。共享底座的适配器可以共用部分计算，但仍须按真实计算图计数。来源本身准确并不免除接口、尺度和组合规则的检查。

\subsubsection{基于条件规则的预测集成}
\label{sec:ket-integration-conditional-prediction}

一个全局权重无法表达“清晰照片相信模型一，模糊照片相信模型二”。可以让选择器观察输入质量、任务标识或可获得的证据，输出条件权重。以$u$标识目标任务$\mathcal T_u$，与来源索引$k$区分；设可观察信息为$\bm{q}(\bm{x},u)$，则
\begin{equation}
  \bar p(y\mid\bm{x},u)
  =\sum_k\alpha_k(\bm{q}(\bm{x},u))p_k(y\mid\bm{x},u).
  \label{eq:ket-integration-conditional-average}
\end{equation}
若结果仍解释为混合分布，权重仍须非负且归一。只有一个权重为一时是硬选择；多个权重非零时是软组合。选择器可以在带质量信息和真实标签的独立材料上学习各来源的条件损失，再决定权重；也可以训练一个以组合损失为目标的门控函数。数据应覆盖拟服务的质量范围，来源在选择器训练样本上的过拟合预测不能当作可靠性证据。

考虑两张教学照片，一张清晰的月季，另一张模糊的非月季。模型一给出的正类概率分别为$0.9,0.95$，模型二给出$0.2,0.2$。始终平均得到$0.55,0.575$，第二张被错判；若可观察的清晰度指标把前者分到清晰区域、后者分到模糊区域，分别调用模型一和模型二便得到$0.9,0.2$，两张均正确。选择器使用的是图像清晰度，不是“哪个模型此时答对”这个预测时未知的量。若清晰度判断错误，两次路由恰好都选错来源，条件组合反而更差。

按任务确定一次权重，与每张输入重新选择来源具有不同成本。前者根据该任务的说明或少量示例求$\bm{\alpha}_u$，随后任务内固定使用；在多任务服务范围内，它仍是条件组合。后者每次都需运行选择器。若清晰度可以低成本从原图计算，硬路由可能只需一个来源；若选择依据是各模型已经算出的置信度，则所有来源早已执行，丢弃其中大部分答案不等于节省来源计算。

当输入超出已验证范围、证据缺失或选择器无法可靠区分来源时，可以按预先确定的规则弃权，转人工复核或更保守的备用流程。弃权阈值也需验证，并报告覆盖率与被接受样本的错误率。地区、季节或设备变化可能改变“清晰度与来源可靠性”的联系；置信度失真、路由误差和条件失配都可能抵消来源互补性。逐输入选择仍只是在已有答案之间分配信任，不能据此声称组合了来源内部的推理步骤。

\subsection{计算模块的组合}
\label{sec:ket-integration-modules}

若目标是“从图像找出红色且花瓣呈钟形的花”，来源给出一个物种标签可能不足以继续计算。系统需要传递属性、位置或其他中间表示，再由后续模块完成判断。此时组合对象从答案集合转为计算图：必须说明每个模块接收什么、返回什么，以及哪些中间结果流向哪些模块。图~\ref{fig:ket-integration-interfaces}将这种接口与同义预测的汇总作了对照。

\begin{figure*}[htbp]
  \centering
  % Interface comparison, two fixed tasks; schematic, no empirical data.
% Canvas: 164 mm x 69 mm. Dashed arrows carry control, not predictions.
\begin{tikzpicture}[
  x=1mm,y=1mm,font=\figurefont\footnotesize,text=BookInk,
  box/.style={rectangle,draw=FigureModel,fill=FigureModel!7,
    line width=0.8pt,minimum width=22mm,minimum height=9mm,
    align=center,inner sep=1mm},
  input/.style={box,draw=FigureInput,fill=FigureInput!7,
    minimum width=15mm},
  result/.style={box,draw=FigureOutput,fill=FigureOutput!7,
    minimum width=20mm},
  control/.style={box,draw=BookAmber,fill=BookAmber!7},
  flow/.style={-{Latex[length=2mm,width=1.3mm]},draw=BookInk,
    line width=0.8pt,rounded corners=1mm},
  choose/.style={flow,dashed,draw=BookAmber},
  edge label/.style={fill=white,inner sep=0.6mm,
    font=\figurefont\footnotesize}
]
  \path[use as bounding box] (0,0) rectangle (164,69);
  \node[anchor=west,font=\figurefont\small] at (0,65) {(a) 同义预测的集成};
  \node[anchor=west,text=BookMuted] at (0,58) {任务：识别物种};
  \node[input] (image-a) at (9,33) {图像};
  \node[box] (model-a) at (39,45) {来源$f_1$};
  \node[box] (model-b) at (39,21) {来源$f_2$};
  \node[result] (aggregate) at (68,33) {聚合};
  \node[result] (species) at (68,6) {物种};
  \draw[flow] (image-a.east) -- (22,33) |- (model-a.west);
  \draw[flow] (22,33) |- (model-b.west);
  \draw[flow] (model-a.east) -- (68,45) -- (aggregate.north);
  \draw[flow] (model-b.east) -- (68,21) -- (aggregate.south);
  \node[edge label] at (65,51) {$p_1(y\mid\bm{x})$};
  \node[edge label] at (64,15) {$p_2(y\mid\bm{x})$};
  \draw[flow] (aggregate.east) -- (80,33) |- (species.east);
  \draw[BookRule,line width=0.6pt] (85,1)--(85,68);

  \node[anchor=west,font=\figurefont\small] at (90,65) {(b) 中间计算的组合};
  \node[anchor=west,text=BookMuted] at (90,58) {任务：属性条件判断};
  \node[input] (image-b) at (101,43) {图像};
  \node[box] (attributes) at (144,43) {属性模块$g$};
  \node[control] (selector) at (104,21) {条件选择器};
  \node[box] (reasoner) at (144,21) {判断模块$h$};
  \node[result] (decision) at (144,5) {判断};
  \draw[flow] (image-b.east) -- (attributes.west);
  \draw[flow] (attributes.south) -- (reasoner.north);
  \node[edge label,anchor=west] at (147,32) {$\bm{z}$};
  \node[edge label,anchor=east] at (142,32) {中间表示};
  \draw[flow] (image-b.south) -- (101,32) -| (selector.north);
  \draw[choose] (selector.east) -- node[edge label,above] {选择信号} (reasoner.west);
  \draw[choose] (selector.west) -- (89,21) -- (89,52)
    -| (attributes.north);
  \draw[flow] (reasoner.south) -- (decision.north);
\end{tikzpicture}

  \caption{预测集成与模块组合的接口差别。左面板固定物种分类任务，两条接口都携带相同类别顺序的预测分布；右面板固定属性判断任务，上游输出中间表示，下游据此继续计算，虚线只携带选择信号。若上游向量依次表示“红色、钟形”，下游却按“钟形、红色”读取，维度虽然相同，判断含义已经改变。能连线只说明形式可接，不说明整条链正确。}
  \label{fig:ket-integration-interfaces}
\end{figure*}

\subsubsection{基于固定连接的模块组合}
\label{sec:ket-integration-fixed-modules}

固定连接在约定范围内使用同一计算图。例如属性模块$g$把图像映射为$\bm{z}=g(\bm{x})\in[0,1]^2$，两维依次是“红色”和“钟形”的置信度；判断模块$h$依据这两个值回答查询。教学规则取$h(\bm{z})=\mathbf{1}\{z_1>0.8,\ z_2>0.6\}$。若$g(\bm{x})=(0.9,0.7)^\top$，链式计算$h(g(\bm{x}))$返回真。若下游误以为两维顺序相反，相同向量会按错误阈值判断而返回假。若接口正确，但上游把第二个属性错估为$0.4$，结果同样为假，原因却是上游误差，而非接口语义错位。

真实接口比这两个数更复杂。除维度、类别顺序和单位，还须匹配归一化尺度、空间坐标、序列位置、缺失值约定和所需状态。一个处理视频的模块可能要求连续帧和历史缓存，不能只凭单帧张量的形状接入。串联形式为$h(g(\bm{x}))$；并联交互可以把颜色与形态模块的输出拼接为$h([g_{\mathrm{color}}(\bm{x});g_{\mathrm{shape}}(\bm{x})])$。拼接只是保留两个接口，是否需要交互项由下游目标决定。

接口不一致时，可在配对样本上训练映射$r_{\bm{\phi}}$，使$h(r_{\bm{\phi}}(g(\bm{x})))$工作。若映射目标是下游熟悉的表示$\bm{z}^{\mathrm{ref}}$，可用$\sum_i\lVert r_{\bm{\phi}}(g(\bm{x}_i))-\bm{z}^{\mathrm{ref}}_i\rVert_2^2$；若只有最终标签，可在固定来源的前提下通过整链损失训练$\bm{\phi}$。前者需要中间对应，后者需要可传梯度或其他可行优化途径，且不保证映射恢复了预设的属性语义。应明确更新的只是接口，还是还包括上下游参数。

下游若只在真实属性标签上训练，训练输入通常接近零或一；上游实际输出的软概率、相关错误和缺失属性会构成新的输入分布。因而应分别用真实中间值和预测中间值测量下游表现，定位接口误差与上游误差，也可用来源预测构造下游训练材料。冻结来源参数只保持其原计算在原输入、状态及调用方式下不变；串联后输入分布和决策路径已经改变，不能由“没有改来源参数”推出组合系统保留全部来源任务的效果。

\subsubsection{基于条件选择的模块组合}
\label{sec:ket-integration-conditional-modules}

多个目标可能需要不同的属性模块和连接方式。为限制搜索空间，可以先规定\term{组合语法}（composition grammar），即允许哪些模块出现在什么位置、哪些输出能够连接到哪些输入。\term{模块化元学习}（modular meta-learning）在多个相关任务上准备可重组模块，再在新任务的少量材料上选择结构；这里关注模块库已形成之后的结构选择。任务分布及留出条件见第~\ref{sec:ket-transfer-task-splits}~节，为新任务准备可学习知识的目标见第~\ref{sec:ket-transfer-preparation}~节\scite{ket-alet2018-modular}

设模块库为$g_1,g_2,h_1,h_2$，前两个是属性提取器，后两个是与其输出格式相容的判断器，候选结构为$\mathcal{G}=\{h_b\circ g_a:a,b\in\{1,2\}\}$。在新任务支持集$S_u$上选
\begin{equation}
  G_u^\star\in\argmin_{G\in\mathcal{G}}
  \frac{1}{|S_u|}\sum_{(\bm{x},y)\in S_u}\ell(G(\bm{x}),y).
  \label{eq:ket-integration-structure-search}
\end{equation}
这是在固定模块参数下搜索计算结构。假设四个候选在四个支持样本上分别错$2,1,3,0$个，零一经验风险便是$0.5,0.25,0.75,0$，选择第四个。独立查询集$Q_u$才用于检查这一选择是否推广到未见样本；支持集上的零错误不是泛化保证。

候选很多时，不必穷举所有结构。模块化元学习采用局部结构变换，例如替换一个模块或增删一个合法连接，并用模拟退火搜索。误差不增加时接受；误差增加时，在同一温度下增量越大，接受概率越低。同一正误差增量下，温度越低，接受概率也越低，搜索中逐步降温便减少了接受劣解的机会。本例也可从$h_1\circ g_1$出发，替换判断器，再替换提取器。有限搜索预算下可能停在次优结构，不能把搜索算法写成最优性保证。若任务内固定$G_u^\star$，预测只运行选中的模块；若每个指令或输入决定结构，还须计入在线路由及结构切换。

条件选择也可以发生在同一层的中间表示上。AdapterFusion冻结已训练的底座和来源适配器，另训练融合层，使当前上下文决定各适配器表示的权重。
设某层的底座表示为$\bm{h}$，来源适配器输出为$\bm{z}_k\in\mathbb{R}^{d}$，可将其核心注意力运算写为
\begin{equation}
  \begin{aligned}
    \bm{q}&=\bm{W}_{Q}\bm{h},&
    \bm{k}_k&=\bm{W}_{K}\bm{z}_k,&
    \bm{v}_k&=\bm{W}_{V}\bm{z}_k,\\
    a_k&=\frac{\exp(\bm{q}^{\top}\bm{k}_k)}
                  {\sum_j\exp(\bm{q}^{\top}\bm{k}_j)},&
    \bm{z}_{\mathrm{fuse}}&=\sum_k a_k\bm{v}_k.
  \end{aligned}
  \label{eq:ket-integration-adapter-attention}
\end{equation}
查询、键和值的投影维度须相容，融合输出再接回后续网络；这里省略层与位置索引及残差结构。目标任务标签用于训练这些投影和相应任务头，来源模块不更新。权重随表示变化，因此是输入条件组合，而不是离线平均适配器参数。第~\ref{sec:ket-sharing-conditional}~节的MMoE还让任务共同训练专家；门控形式相近，不代表共享量的学习条件相同。软注意力一般仍需计算所有候选适配器，不能自动获得稀疏路由的计算节省。原论文主要实验还包含目标任务自身已训练的适配器，这与完全未见目标任务的模块重组条件不同\scite{ket-pfeiffer2021adapterfusion}

判断是否产生组合泛化，需要把来源成分与组合关系分开留出。属性实验可以让“红色”和“钟形”各自出现，却把“红色且钟形”组合留作测试；指令实验可以让“筛选红花”和“按高度排序”分别出现，测试“筛选红花后排序”。还要比较单一来源选择、固定组合和真正重组，检查系统是否利用了中间结果。为组合泛化构造元学习任务的研究，用相似性驱动的采样形成训练与推广任务对，使相同成分出现在不同组合中；它改变了来源经验的安排，不是证明任意可连接模块都能泛化\scite{ket-conklin2021compositional}

技能重组还涉及动作执行后的状态，不能照搬静态属性接口。花卉采样机器人若已有“接近植株”和“对准相机”技能，前一技能终止时的状态必须落在后一技能能处理的范围。\term{后继特征}（successor features）记录既定策略下未来特征的折扣累计期望。在共享动力学、特征表示和状态动作接口，并用特征的线性函数表达奖励时，把后继特征与新奖励权重相乘，可以评价已有策略在新任务上的价值，再以广义策略改进综合多个策略的价值信息。这不是任意技能串接保证；完整决策、探索与价值推导属于第~\ref{part:reinforcement-learning}~部分的内容范围\scite{ket-barreto2016successor}

长期模块库还需记录每个模块的版本、能力范围、接口和依赖。重复能力使搜索空间增长，淘汰模块可能破坏依赖它的组合，检索遗漏则使可用模块没有进入候选。应保留可复现的组合清单和回退版本，分别核算检索、搜索、加载及执行成本。版本与依赖管理见第~\ref{sec:ket-update-versions}~节，载体及访问条件见第~\ref{sec:ket-overview-information}~节；这里要检验的是已有能力能否按约定共同调用。

\section{基于参数融合的知识整合}
\label{sec:ket-integration-merging}

若花卉终端只能容纳一份模型，保留多个分类器并逐一运行可能超出预算。参数融合尝试直接生成一份承接来源能力的参数，部署时不再逐个查询来源。但参数数组并不是标有语义名称的能力清单：相同位置可能承担不同计算，相近的权重也可能处在不同的决策边界两侧。本节依次确定参数对应、构造融合规则，再利用允许的信息选择系数；没有对应关系，后两步的算术就可能失去意义。

\subsection{参数对应关系的建立}
\label{sec:ket-integration-alignment}

同架构只保证参数形状可比，并不保证坐标语义相同。考虑具有逐元素激活$\sigma$的双层网络
\[
  f(\bm{x})=\bm{W}_2\sigma(\bm{W}_1\bm{x}+\bm{b}_1)+\bm{b}_2,
  \quad \bm{W}_1\in\mathbb{R}^{m\times d},\
  \bm{W}_2\in\mathbb{R}^{c\times m}.
\]
设$\bm{P}\in\mathbb{R}^{m\times m}$为置换矩阵，只改变隐藏单元的排列。由于$\bm{P}^{\top}\bm{P}=\bm{I}$且$\sigma(\bm{P}\bm{v})=\bm{P}\sigma(\bm{v})$，配套变换
\begin{equation}
  \begin{aligned}
    \widetilde{\bm{W}}_1&=\bm{P}\bm{W}_1,&
    \widetilde{\bm{b}}_1&=\bm{P}\bm{b}_1,\\
    \widetilde{\bm{W}}_2&=\bm{W}_2\bm{P}^{\top},&
    \widetilde{\bm{b}}_2&=\bm{b}_2
  \end{aligned}
  \label{eq:ket-integration-permutation}
\end{equation}
不改变函数：新网络输出为$\bm{W}_2\bm{P}^{\top}\bm{P}\sigma(\bm{W}_1\bm{x}+\bm{b}_1)+\bm{b}_2=f(\bm{x})$。前一层换行时，后一层必须同步换列；只改一层通常不再等价。

\begin{example}[完全相同的函数也不能直接平均]
  \label{ex:ket-integration-permutation}
  令输入、输出均为标量，$\sigma=\operatorname{ReLU}$，偏置为零，
  \[
    \bm{W}_1=\begin{bmatrix}1\\-1\end{bmatrix},\quad
    \bm{W}_2=\begin{bmatrix}1&1\end{bmatrix},\quad
    \bm{P}=\begin{bmatrix}0&1\\1&0\end{bmatrix}.
  \]
  原网络与置换后的网络都计算$\operatorname{ReLU}(x)+\operatorname{ReLU}(-x)=|x|$。直接逐坐标平均两份权重，第一层却变为零向量，输出恒为零。在$x=1$时两个来源都输出$1$，融合模型输出$0$。若先用逆置换将第二个来源对齐，两个参数数组恢复一致，平均才仍计算$|x|$。
\end{example}

Git Re-Basin把这类隐藏单元对应作为显式问题：可以在共同输入上比较激活并求线性指派，也可以只用权重，通过逐层更新置换来增大匹配程度。每次改变一层置换，都须保持邻层和残差等结构所要求的配套关系。权重匹配不要求读取训练样本，激活匹配则需要代表性输入；二者获得的信息不同。搜索到的对应也不保证所有插值点都保留低损失。该研究还人为构造了无法通过任何单元置换实现线性连通的非SGD解，说明连通性不能推广到任意解；这一反例并未反驳对SGD所得解的限定猜想\scite{ket-ainsworth2023rebasin}

两个来源若分别学习物种和病害，强制所有层一一对应还可能抹去各自专用的特征。ZipIt!在同架构模型上利用无标签输入的激活相关性，允许跨模型以及同一模型内部的冗余特征合并，并允许只融合到某层，保留后续任务分支。它不只是给模型二找一个置换：允许模型内部匹配，是为了避免把不相似的跨任务特征硬配成一对。保留分支意味着部署仍可能需要任务身份以及对应输出头，也未获得整个模型的单一路径成本\scite{ket-stoica2024zipit}

低秩模块还存在另一层坐标问题。\term{低秩适应}（low-rank adaptation，LoRA）把权重更新写成两个瘦矩阵的乘积，仅学习或保存相对底座的低秩增量。对某一层，记$\Delta\bm{W}_k=\bm{B}_k\bm{A}_k$，其中$\bm{B}_k\in\mathbb{R}^{d_{\mathrm{out}}\times r}$、$\bm{A}_k\in\mathbb{R}^{r\times d_{\mathrm{in}}}$；固定缩放量已吸收入因子，其训练位置和参数开销见第~\ref{sec:ket-transfer-added-parameters}~节。即使增量相同，对任意可逆$\bm{R}$也有$(\bm{B}_k\bm{R})(\bm{R}^{-1}\bm{A}_k)=\Delta\bm{W}_k$，所以因子的内部坐标不唯一。不同任务的更新还可能作用于不同子空间；因子维度相同不等于其基向量相同。

KnOTS先形成完整增量矩阵，再逐层拼接并作联合奇异值分解，使各任务增量使用一个共同的左侧基。例如令
\begin{equation}
  \begin{aligned}
    \relax[\Delta\bm{W}_1\ \cdots\ \Delta\bm{W}_K]
      &=\bm{U}\bm{\Sigma}[\bm{V}_1^\top\ \cdots\ \bm{V}_K^\top],\\
    \Delta\bm{W}_k&=\bm{U}\bm{\Sigma}\bm{V}_k^\top.
  \end{aligned}
  \label{eq:ket-integration-knots}
\end{equation}
其中$\bm{U}\in\mathbb{R}^{d_{\mathrm{out}}\times q}$、$\bm{\Sigma}\in\mathbb{R}^{q\times q}$、$\bm{V}_k^\top\in\mathbb{R}^{q\times d_{\mathrm{in}}}$，$q$为拼接矩阵的秩，不截断时等式精确成立。随后在共同坐标中的$\bm{V}_k^\top$上应用融合规则，再乘回$\bm{U}\bm{\Sigma}$。这与融合完成后为压缩结果再作低秩分解不同。特别地，纯线性求和在精确分解下与先求和再变换等价；改变坐标主要会改变裁剪、符号选择等依赖坐标的非线性操作，不能把分解本身写成凭空增加能力\scite{ket-stoica2025knots}

保持共同底座时，来源增量至少引用相同的层、输入输出坐标和基准参数；若只保存适配器，还必须保留底座的准确版本。重新训练的独立底座没有这种直接对应，不能将其适配器原样相加后装到任意一个底座上。此时通常需要访问完整底座、建立层与激活对应，或改用可查询行为的蒸馏。异构深度、宽度和算子还可能使参数运算根本未定义，不能靠补零或同名层匹配便认定兼容。共同起点是有利条件，不是融合成功的充分条件。

\subsection{融合规则的构造}
\label{sec:ket-integration-merging-rule}

对应明确之后，仍需约定参与运算的是全部参数、共享层，还是相对底座的更新。记来源参数为$\bm{\theta}_k\in\mathbb{R}^{d_\theta}$，底座为$\bm{\theta}_0$。线性完整参数组合与增量组合的联系是
\begin{equation}
  \sum_k\alpha_k\bm{\theta}_k
  =\left(\sum_k\alpha_k\right)\bm{\theta}_0+
    \sum_k\alpha_k(\bm{\theta}_k-\bm{\theta}_0).
  \label{eq:ket-integration-rule-equivalence}
\end{equation}
当$\sum_k\alpha_k=1$时，完整参数平均与加回底座的同系数增量组合相同；不满足这一条件时，底座的系数也发生变化。若另做裁剪或因子相乘，则未必还能用这一线性恒等式改写。本小节给出规则与待定系数，具体数值由下一小节确定。

\subsubsection{完整参数的融合}
\label{sec:ket-integration-full-weights}

Model Soups从同一个预训练模型出发，在同一目标任务上用不同学习率、数据增强或训练随机性得到多个微调结果，然后平均其权重。均匀方案用$\alpha_k=1/K$；贪心方案先按验证成绩排序，从最佳来源开始，依次尝试加入新来源，只有候选平均在验证集上不下降时才保留；学习型方案还可以优化插值系数。共同微调起点使坐标更可能对应，但不合适的来源仍可能破坏平均结果，因此验证接纳规则是方法的一部分\scite{ket-wortsman2022-soups}

等权平均把每个参数坐标视为同等可靠，数据感知方法则利用来源训练统计改变逐坐标或逐层权重。Fisher merging以各来源参数附近的对角Fisher信息近似局部曲率，对坐标$i$可写成
\[
  \theta_i^{\mathrm{F}}
  =\frac{\sum_k F_{k,i}\theta_{k,i}}
         {\sum_k F_{k,i}},
\]
其中分母为零的坐标须采用预定回退规则。它需要来源数据或已保存的Fisher统计，且对角近似忽略参数相关；“重要”是相对于来源似然的局部量，不保证多任务融合后最优\scite{ket-matena2022-fisher}

RegMean针对线性层保留输入二阶统计。若第$k$个来源层为$\bm W_k$，输入Gram矩阵为$\bm G_k=\bm X_k^\top\bm X_k$，最小化各来源层输出的平方偏差得到
\[
  \bm W^{\mathrm{R}}
  =\left(\sum_k\bm G_k\right)^{-1}
    \left(\sum_k\bm G_k\bm W_k\right),
\]
实际实现需处理转置约定、正则化与病态矩阵。该路线不要求保存原始文本标签，却需要与来源分布相符的输入统计；统计失配或层坐标不对应时，闭式解仍可能融合失败\scite{ket-jin2023-regmean}

参数平均与预测平均保留的信息不同。前者只执行$f(\bm{x};\sum_k\alpha_k\bm{\theta}_k)$，后者执行所有$f(\bm{x};\bm{\theta}_k)$后才汇总。只有模型对所融合参数具有相应线性结构时，两种顺序才可能交换；多层乘法和激活函数一般破坏这种性质。

\begin{example}[非线性使两种平均不相等]
  \label{ex:ket-integration-nonlinear-average}
  为隔离激活的影响，取单参数分数模型$f(x;w)=\operatorname{ReLU}(wx)$，两个来源为$w_1=2,w_2=-1$，固定$x=1$及相等权重。平均来源分数得到$(2+0)/2=1$，平均参数却是$\bar w=0.5$，其分数$f(1;\bar w)=0.5$。

  若用二分类分数向量$(f(x;w),0)$生成概率，两个来源的正类概率分别为$e^2/(1+e^2)\approx0.8808$和$0.5$，平均概率约为$0.6904$；融合模型概率为$e^{0.5}/(1+e^{0.5})\approx0.6225$。即使改为先平均分数再做softmax，所得$e^1/(1+e^1)\approx0.7311$也不同。来源、输入和权重都未改变，改变的只是平均发生的位置。
\end{example}

同一花卉分类任务的多个解，至少共享输出含义和主要评价目标；将物种识别与病害识别的能力合到一起，还要保留两个任务不同的目标和接口。前一种设定下的实证成功不能直接推广到任意技能融合。若完整输出头的类别含义不同，应只融合可对应的共享层，保留独立任务头；使用哪个头可能需要任务身份，未知任务的选择问题并未因此消失。

融合验证应分别检查每个来源任务的风险，不能只看平均成绩。对同一输入只运行共享主干和一个头，可能比两个完整来源便宜；若需要执行多个分支、重新归一化统计量或额外校准，成本应按实际结构计算。参数之外的运行状态也要处理，例如批归一化的统计量不能在未检查含义时一并当作普通权重平均。

来源训练目标也可以改变。Self-Soupervision研究共同底座在不同自监督目标或数据上继续训练后如何形成可融合来源；部分实验随后仍用标签微调，另有直接融合自监督来源并按少量标签选择的方案，以及无监督选择的变体。因此，“来源的继续训练不用标签”“融合系数的确定不用标签”和“最终评价不用标签”是不同陈述。接触测试分布的无标签输入时，还须按直推式（亦称转导式）或适应协议报告，不能当作未见测试数据的普通归纳预测\scite{ket-fuller2026selfsoup}

\subsubsection{参数增量的融合}
\label{sec:ket-integration-deltas}

若希望保留共同底座，把来源差异显式写出来更方便。\term{任务向量}（task vector）是某个微调结果相对同一底座的参数差，并不表示已经分离出一项纯粹的语义技能。任务算术依次取差、缩放组合、加回底座：
\begin{equation}
  \bm{\tau}_k=\bm{\theta}_k-\bm{\theta}_0,\qquad
  \bm{\theta}_{\mathrm{merge}}(\bm{\lambda})
    =\bm{\theta}_0+\sum_{k=1}^{K}\lambda_k\bm{\tau}_k.
  \label{eq:ket-integration-task-arithmetic}
\end{equation}
常见的原始多任务形式让所有$\lambda_k$共享一个缩放系数，随后可扩展为来源各自的系数。与概率混合不同，参数系数不必非负或和为一，但扩大候选范围也会增加失效风险\scite{ket-ilharco2022taskarithmetic}

TIES-Merging针对增量中的小幅变化和逐坐标符号冲突，依次执行裁剪、符号选举和同号合并。令$\widetilde{\bm{\tau}}_k$保留$\bm{\tau}_k$中绝对值较大的给定比例坐标，其余置零；对每个坐标$j$，按带幅度的和选择方向，并只平均与该方向一致的非零值：
\begin{equation}
  \begin{aligned}
    s_j&=\operatorname{sign}\!\left(\sum_k\widetilde{\tau}_{k,j}\right),\\
    I_j&=\{k:\widetilde{\tau}_{k,j}\neq0,\
                    \operatorname{sign}(\widetilde{\tau}_{k,j})=s_j\},\\
    \tau_{\mathrm{TIES},j}
      &=\frac{1}{|I_j|}\sum_{k\in I_j}\widetilde{\tau}_{k,j},\\
    \bm{\theta}_{\mathrm{merge}}
      &=\bm{\theta}_0+\lambda\bm{\tau}_{\mathrm{TIES}}.
  \end{aligned}
  \label{eq:ket-integration-ties}
\end{equation}
为使此处算法完整，约定$s_j=0$或$I_j$为空时该坐标输出零，幅度裁剪遇到并列值按固定索引处理。实际实现应同时记录这类退化约定、保留比例及最终缩放。符号选举并非简单按正负来源个数投票，幅度也参与决定；同号平均的分母是参与该坐标的来源数，不总是$K$\scite{ket-yadav2023ties}

取教学底座及两个来源为
\[
  \begin{aligned}
    \bm{\theta}_0&=(1,1,1)^\top,\\
    \bm{\theta}_1&=(1.4,0.8,1.01)^\top,\\
    \bm{\theta}_2&=(1.2,1.3,0.98)^\top.
  \end{aligned}
\]
逐坐标相减得到
\[
  \begin{aligned}
    \bm{\tau}_1&=(0.4,-0.2,0.01)^\top,\\
    \bm{\tau}_2&=(0.2,0.3,-0.02)^\top.
  \end{aligned}
\]
每个保留两个最大幅度坐标，得到$(0.4,-0.2,0)^\top$和$(0.2,0.3,0)^\top$。前两维选举方向均为正，合并增量为$(0.3,0.3,0)^\top$：第一维平均两个来源，第二维只保留$0.3$。原始求和则是$(0.6,0.1,-0.01)^\top$。缩放均取一，加回底座后，TIES得到$(1.3,1.3,1)^\top$，任务算术得到$(1.6,1.1,0.99)^\top$。这个算例展示了算术差别，不声称哪个向量对应更好的任务表现。

DARE改变的主要是增量稀疏化方式：以概率$p$随机丢弃坐标，再把保留值除以$1-p$，之后仍可接任务算术或TIES等规则。用$M_{k,j}\sim\operatorname{Bernoulli}(1-p)$表示保留掩码，$0\leq p<1$，有
\begin{equation}
  \begin{aligned}
    \widehat{\tau}_{k,j}&=\frac{M_{k,j}}{1-p}\tau_{k,j},\\
    \mathbb{E}_M[\widehat{\tau}_{k,j}]&=\tau_{k,j},&
    \operatorname{Var}_M(\widehat{\tau}_{k,j})
      &=\frac{p}{1-p}\tau_{k,j}^{2}.
  \end{aligned}
  \label{eq:ket-integration-dare}
\end{equation}
重缩放保持的是对随机掩码的坐标期望；固定输入的单层线性更新也因期望的线性性保持其期望输出，但非线性网络输出及任务损失不因此无偏。当$p$接近一，单坐标方差反而增大。该方法的实证依赖增量的尺度、冗余与丢弃率，不能把对小幅微调增量的观察推广到任意独立训练模型\scite{ket-yu2024dare}

例如$p=0.5$，对前述$\bm{\tau}_1$取保留掩码$(1,0,1)^\top$，得到$(0.8,0,0.02)^\top$；它不是原增量，只是在所有掩码上的期望等于原增量。TIES按幅度和符号有选择地去掉坐标，DARE随机去掉坐标再补偿期望，二者可以组合，却不能互称为同一个裁剪步骤。

参数符号取决于坐标，正负冲突不等于语义冲突。隐藏单元重排甚至可以在函数不变时改变冲突统计。类似地，$\bm{\theta}_0-\lambda\bm{\tau}_k$只是沿取反的参数方向移动，不是把某项知识从系统中精确删除；若要满足训练影响消除要求，仍须采用第~\ref{sec:ket-update-unlearning-validation}~节及第~\ref{chap:trust-privacy}~章的目标与证据。

LoRAHub则直接对低秩因子进行加权。在与前文一致的$\Delta\bm{W}_k=\bm{B}_k\bm{A}_k$记号下，对相同底座、插入层和秩的来源构造
\begin{equation}
  \begin{aligned}
    \bm{B}(\bm{w})&=\sum_k w_k\bm{B}_k,&
    \bm{A}(\bm{w})&=\sum_k w_k\bm{A}_k,\\
    \Delta\bm{W}(\bm{w})&=\bm{B}(\bm{w})\bm{A}(\bm{w}).
  \end{aligned}
  \label{eq:ket-integration-lorahub}
\end{equation}
系数可正可负，不要求归一。只取两个来源时，乘积展开为
\begin{equation}
  \begin{aligned}
    \Delta\bm{W}(\bm{w})
      ={}&w_1^2\bm{B}_1\bm{A}_1+w_2^2\bm{B}_2\bm{A}_2\\
       &+w_1w_2(\bm{B}_1\bm{A}_2+\bm{B}_2\bm{A}_1).
  \end{aligned}
  \label{eq:ket-integration-lorahub-cross-terms}
\end{equation}
来源自身的项带系数平方，还出现跨来源因子的乘积，所以一般不等于$\sum_k w_k\Delta\bm{W}_k$。LoRAHub原文把输出侧因子命名为$\bm{A}$，这里为与本章统一而交换因子字母，运算次序与维度不变\scite{ket-huang2024lorahub}

取秩一教学因子$\bm{B}_1=(1,0)^\top$、$\bm{A}_1=(1,0)$、$\bm{B}_2=(0,1)^\top$、$\bm{A}_2=(0,1)$及$w_1=w_2=1/2$。直接平均完整增量得到$\operatorname{diag}(1/2,1/2)$，秩为二；因子组合得到每个元素均为$1/4$的$2\times2$矩阵，秩为一。非对角元素来自交叉项，差别既包括数值也包括可表达的秩。完成组合后只运行形成的低秩模块，必要时将增量加进底座权重；这不同于保留两个适配器、逐输入混合其激活的AdapterFusion。

\subsection{融合系数的确定}
\label{sec:ket-integration-coefficients}

融合公式只定义候选模型，不决定哪个候选值得部署。固定系数可以来自预先约定，例如等权平均；有标签验证选择可直接最小化各目标任务的经验风险；无标签适应只能利用熵、稳定性等替代量。无论选哪一种，都应先声明候选范围、信息访问和计算预算。若反复观察测试标签再选系数，报告的提高就混入了选择收益。

LoRAHub针对一个新任务取得少量输入输出示例，固定底座和来源因子，把$\bm{w}$作为低维搜索变量。每次提议一组系数，按式~\eqref{eq:ket-integration-lorahub}生成候选模块，在这些示例上计算目标输出的负对数似然，并加入系数幅度正则；无梯度优化器据此提出下一候选。达到评估预算后选定系数，任务内不再逐样本重搜。

原论文用Shiwa选择无梯度优化策略，多数实验采用协方差矩阵自适应进化策略，并限制候选模块规模。“无梯度”省去了搜索中的反向传播，目标示例、各候选的前向评分和搜索预算仍然不可少。

可以把上述搜索目标概括为
\begin{equation}
  \bm{w}_u^\star\in\argmin_{\bm{w}\in\Omega}
   \left\{-\frac{1}{|S_u|}
   \sum_{(\bm{x},y)\in S_u}\log p_{\bm{w}}(y\mid\bm{x})
   +\eta\lVert\bm{w}\rVert_1\right\}.
  \label{eq:ket-integration-lorahub-search}
\end{equation}
这里$\Omega$为预定搜索范围，$\eta$控制正则强度；序列输出的对数概率由对应词元条件概率相加。$S_u$承担适应作用，最终评价仍需独立材料。增加候选数会提高搜索维度，少量示例也可能让选择过拟合；不能因为只学习几个系数就免除这些问题。

AdaMerging不用目标标签，而用无标签输入上的预测熵优化任务向量系数。任务级形式为$\bm{\theta}=\bm{\theta}_0+\sum_k\lambda_k\bm{\tau}_k$；层级形式给每层$\ell$独立系数，
\begin{equation}
  \begin{aligned}
    \bm{\theta}_{\ell}(\bm{\lambda})
      &=\bm{\theta}_{0,\ell}+\sum_k\lambda_{k,\ell}\bm{\tau}_{k,\ell},\\
    \bm{\lambda}^{\star}
      &\in\argmin_{\bm{\lambda}}\frac{1}{|U|}
        \sum_{\bm{x}\in U}H(p_{\bm{\lambda}}(\cdot\mid\bm{x})).
  \end{aligned}
  \label{eq:ket-integration-adamerging}
\end{equation}
$U$表示无标签适应输入集，$H(p)=-\sum_y p(y)\log p(y)$。从预设系数初始化，固定底座与任务向量，通过候选模型的前向计算和熵梯度更新系数，直到预算耗尽或预定停止条件满足。原研究在多任务无标签测试样本上进行此过程，并按任务接口计算输出；任务级系数意指每个来源向量一个系数，不等于每个输入重新生成一组系数\scite{ket-yang2024adamerging}

低熵只说明预测集中。恒将每张花卉照片判成月季也可获得接近零的熵，却可能大量出错。因此需要检查类别坍缩、来源任务损失和分布变化，不能把熵下降当成正确率提高的证明。若$U$来自当前测试批次，应明确这是直推式预测或测试时适应；若声称处理未来独立输入，应冻结系数后另测。适应轮数、可见任务身份、是否持续更新以及重置规则都属于协议。

有些来源能力无法在一份给定容量的参数中同时达到各自最优。可把目标写成风险向量$(R_1(\bm{\theta}),\ldots,R_K(\bm{\theta}))$，用偏好权重表达哪些损失更不能接受。\term{帕累托最优}（Pareto optimality）表示无法在不损害其他目标的前提下继续改进某一目标：不存在另一个可行模型在所有任务上不差且至少一项更好。该性质相对于给定可行范围而言，实际算法通常只能近似寻找这种折中。

Pareto Merging将偏好感知融合写成多目标问题，并用偏好无关的共享参数与偏好相关的低秩修正构造一族候选。它分别讨论无数据条件下的参数距离替代目标，以及使用无标签数据时的分任务熵目标；这些代理目标的折中不应直接等同于真实风险的全局帕累托最优。对花卉系统，可以分别报告物种与病害任务的性能曲线，再按用途选定候选，而不是用一个平均分掩盖病害漏检的退化\scite{ket-chen2025pareto}

这里在已有来源限定的融合族中寻找折中。第~\ref{sec:ket-sharing-directions}~节则在共同训练过程中协调各任务的下降方向；两者都可用多目标语言描述，但候选空间、训练证据与资源账目不同。

表~\ref{tab:ket-integration-merging-operations}用“运算对象×所需信息”组织融合规则。参数增量与低秩因子是实现层级不同但可组合的对象，数据感知加权则可以作用于完整参数或选定层；它们不是一棵互斥分类树。对应调整是这些规则可以采用的前置步骤：独立初始化、不同任务或低秩子空间不一致时，可按第~\ref{sec:ket-integration-alignment}节寻找置换、激活对应或共同更新子空间，再选择适用规则。

搜索时所需的来源存储、额外校准材料、多次候选融合与前向计算应计入总成本；部署只保存一份参数不代表整合过程免费。连续吸收新来源还可能积累损失，并受到融合顺序影响，因而应保留底座、来源与系数的版本记录，逐任务回归检查，失败时退回已验证版本。是否保有未来可塑性、是否获得未见组合能力，仍须另设学习与组合任务检验。

独立比较还表明，随机挑选任务组合或只使用相近领域会高估融合的组合泛化；来源任务更异质、任务数增加或目标组合被严格留出时，不同方法的排序和失败率都会改变。Tam等的现实化评测属于预印本证据，价值在于提供与方法提出论文不同的压力条件，而不是给出跨模型的最终排名\scite{ket-tam2024-realistic-merging}因此，本章把最佳单来源、预测集成和匹配预算的联合训练保留为必要对照，不能只在能够融合成功的来源集合上报告结果。

\begin{table*}[htbp]
  \centering
  \small
  \setlength{\tabcolsep}{4pt}
  \renewcommand{\arraystretch}{1.2}
  \begin{tabularx}{\linewidth}{@{}>{\raggedright\arraybackslash}p{24mm}
      >{\raggedright\arraybackslash}X
      >{\raggedright\arraybackslash}X
      >{\raggedright\arraybackslash}X@{}}
    \toprule
    操作 & 运算对象与可执行条件 & 系数所用信息 & 部署与主要风险 \\
    \midrule
    完整参数平均 & 全部参数或共享层；层、坐标及输出含义相容，常取共同底座的微调解 & 预设系数，或用验证材料选择、学习系数 & 单模型，可保留多头；可能跨越高损失区域或发生头语义冲突 \\
    \addlinespace
    数据感知加权 & 完整参数或线性层；需保存Fisher或输入二阶统计 & 来源数据、统计缓存或可代表来源的输入；Fisher merging、RegMean & 单模型；统计失配、近似曲率与病态矩阵会放大误差 \\
    \addlinespace
    任务增量融合 & 相对同一准确底座的参数差；坐标相同，可裁剪、选符号或随机稀疏 & 预设系数、标签验证或无标签适应 & 底座加融合增量；可能出现任务干扰、尺度过大或丢弃方差过高 \\
    \addlinespace
    低秩因子组合 & 分别加权因子后相乘；同底座、插入位置和秩，因子形状与约定相容 & LoRAHub按任务用少量示例搜索 & 一个整合模块或写回底座；交叉项、因子坐标不稳及少样本过拟合 \\
    \midrule
    对应调整（前置步骤） & 用置换、激活对应或共同子空间调整参数；异构结构仍可能不兼容 & 权重匹配可不读输入，激活匹配需要输入；系数依后接规则 & 可形成单模型或保留分支；匹配失真及对齐后的任务冲突仍需检查 \\
    \bottomrule
  \end{tabularx}
  \caption{参数融合的条件与操作。上部交叉比较运算对象和所需信息，下部列出可与它们结合的对应调整；同一方法可以同时属于增量、低秩与数据感知路线。无标签只限定监督形式，不代表没有输入数据、任务身份或适应计算；形成单份参数也不保证所有来源能力同时保留。}
  \label{tab:ket-integration-merging-operations}
\end{table*}

\section{基于知识蒸馏的知识整合}
\label{sec:ket-integration-distillation}

两个花卉分类器即使网络结构不同，也可能对同一张照片给出可比较的类别分布。这时无法直接平均参数，却可以把这些分布当作训练约束，让一个承接模型学习。通常把提供信号的来源称为教师，把接受训练的承接模型称为学生。学生不必更小，也不必与教师同架构；关键条件是存在可比较的监督信号、足够的输入覆盖，以及能表达目标行为的学生模型族。教师还可以是本章前面形成的组合系统或融合模型。

\subsection{教师与蒸馏输入的准备}
\label{sec:ket-integration-teachers-inputs}

教师准备从访问条件开始。若只能调用分类接口，就只能使用返回的类别、概率或分数，不能假定能够读取全部隐藏层；若能访问权重和内部激活，才可考虑表示蒸馏。两个教师若对同一输入提供物种概率，需要按第~\ref{sec:ket-integration-prediction}~节检查类别语义、聚合规则和校准；若一个提供物种、另一个提供病害，应保留任务条件或分别设置监督头，不能把两个目标平均成一个不明含义的向量。

教师数量不等于独立知识数量。两个来源可能继承相同训练偏差，也可能在同一输入上给出冲突判断。应明确学生模仿的是等权教师集成、经过校准的条件集成、某个任务专用教师，还是融合后的单模型。若先集成获得更好的监督，再蒸馏给学生，应同时比较单教师、教师集成和学生，不能把集成本身带来的提高全部归因于压缩。来源是否值得信任仍需第~\ref{chap:ket-reuse-selection}~章的判断。

蒸馏输入决定教师哪些行为有机会进入学生。已有训练输入通常贴近教师熟悉的区域，但可能带有记忆性；另采未标注输入能覆盖新的条件，却可能超出教师可靠范围；合成输入可以补充难以采集的样本，也可能偏离真实分布或反复产生相似模板。假设查询集里只有常见月季和百合，没有少见兰花，学生在查询集上的蒸馏损失即使为零，也可以在兰花上任意出错。监督只约束实际查询到的区域，不能从局部拟合推出全部能力保留。

跨模态时还要区分训练与部署观测。广义蒸馏允许教师在训练时使用学生部署阶段不可用的特权信息，例如教师同时看到花卉照片与检测报告，学生部署时只有照片；配对的输入使教师能够为学生可见的照片提供监督\scite{ket-lopezpaz2016generalized}
这需要明确对应哪一株植物、哪一观测时刻，配对表示的具体约束见式~\eqref{eq:ket-update-crossmodal}。若两株植物在学生照片上完全不可区分，却只有检测报告揭示了不同病原，学生不可能逐例恢复仅教师可见的差异。

这一边界也可从期望交叉熵看出。设学生观测为$\bm{x}$，教师额外观测为$\bm{x}^{+}$，软目标为$\bm{q}(\bm{x},\bm{x}^{+})$。不限制学生容量时，对固定$\bm{x}$最小化条件期望交叉熵所得最优预测是
\begin{equation}
  p^\star(y\mid\bm{x})
   =\mathbb{E}\!\left[q_y(\bm{x},\bm{x}^{+})\mid\bm{x}\right].
  \label{eq:ket-integration-privileged-limit}
\end{equation}
因为目标等于$-\sum_y\mathbb{E}[q_y\mid\bm{x}]\log p(y\mid\bm{x})$，其最小值在两分布相等时取得。学生学习的是可见信息条件下的平均教师判断，而不是每个样本被隐藏的额外观测。有限容量、样本不足和教师误差只会增加实际困难。

\subsection{蒸馏约束的构造}
\label{sec:ket-integration-distillation-losses}

教师可以告诉学生最终应该输出什么，也可以提供某层应该形成怎样的表示。关系矩阵、注意力图和教师梯度等信号，都是由一个或多个表示、输出或损失对输入的导数派生出的监督：它们扩大了可比较的结构，却没有形成与“输出/表示”并列且互斥的第三种知识载体。选择信号取决于访问权限、接口可比性和学生结构，多种约束可以联合使用。

\subsubsection{基于输出行为的蒸馏}
\label{sec:ket-integration-output-distillation}

硬类别只保留教师最倾向的答案，忽略其他类别之间的相对关系。对同一张花卉照片，教师认为它更像百合还是菊花，这种次要倾向也可能帮助学生学习决策边界。经典知识蒸馏通过提高softmax温度，使原本很小的概率更容易进入训练信号\scite{hinton2015distilling}

设教师和学生的类别分数分别为$\bm{a}^{\mathrm{T}},\bm{a}^{\mathrm{S}}\in\mathbb{R}^{c}$。温度$T>0$下的分布为
\begin{equation}
  q_j^{(T)}=\frac{\exp(a_j^{\mathrm{T}}/T)}
                    {\sum_{v=1}^{c}\exp(a_v^{\mathrm{T}}/T)},\qquad
  p_j^{(T)}=\frac{\exp(a_j^{\mathrm{S}}/T)}
                    {\sum_{v=1}^{c}\exp(a_v^{\mathrm{S}}/T)}.
  \label{eq:ket-integration-temperature}
\end{equation}
此处上标$(T)$仅表示温度条件，不是学习迭代。$T=1$恢复通常的softmax；在分数不全相等时，升高温度使分布更平缓。以教师分布为目标，可最小化
\begin{equation}
  L_{\mathrm{out}}
  =-\sum_j q_j^{(T)}\log p_j^{(T)},\qquad
  D_{\mathrm{KL}}(\bm{q}^{(T)}\Vert\bm{p}^{(T)})
  =L_{\mathrm{out}}-H(\bm{q}^{(T)}).
  \label{eq:ket-integration-output-loss}
\end{equation}
教师冻结时，其熵与学生参数无关，所以这两个目标对学生具有相同梯度。多教师时，先聚合分数再软化，与分别软化概率再聚合并不等价，应固定一种定义。若接口只返回概率，也须明确采用概率目标或怎样重新温度化，不能假定已经取得原始分数。

\Needspace{10\baselineskip}
\begin{example}[三个类别携带的信息]
  \label{ex:ket-integration-soft-labels}
  对一张教学照片，按“月季、百合、菊花”排序的教师分数为$(2,1,0)$。两个温度下的概率分别约为
  \[
    \begin{aligned}
      T=1:\quad &(0.6652,0.2447,0.0900),\\
      T=2:\quad &(0.5065,0.3072,0.1863).
    \end{aligned}
  \]
  硬目标$(1,0,0)$只保留月季；软目标还说明百合比菊花更有可能。若另一个教师给出$(2,0,1)$，硬类别相同，但次要类别的顺序相反，学生收到的约束也应不同。
\end{example}

温度不仅改变目标的相对倾向，也改变梯度尺度。对学生分数$a_v^{\mathrm{S}}$，
\[
  \frac{\partial\log p_j^{(T)}}{\partial a_v^{\mathrm{S}}}
    =\frac{\mathbf{1}\{j=v\}-p_v^{(T)}}{T}.
\]
把它代入交叉熵，利用$\sum_jq_j^{(T)}=1$，得到
\begin{equation}
  \frac{\partial L_{\mathrm{out}}}{\partial a_v^{\mathrm{S}}}
   =-\frac{1}{T}\sum_jq_j^{(T)}
      \bigl(\mathbf{1}\{j=v\}-p_v^{(T)}\bigr)
   =\frac{p_v^{(T)}-q_v^{(T)}}{T}.
  \label{eq:ket-integration-temperature-gradient}
\end{equation}
这里已经出现一个$1/T$，另一个尺度因子来自概率差本身。对固定的有限分数，令$\bar a^{\mathrm{S}}$和$\bar a^{\mathrm{T}}$为各自类别分数的均值；当温度相对中心化分数足够高时，指数的一阶展开给出
\begin{equation}
  \begin{aligned}
    p_v^{(T)}
      &=\frac{1}{c}+\frac{a_v^{\mathrm{S}}-\bar a^{\mathrm{S}}}{cT}
          +O(T^{-2}),\\
    \frac{\partial L_{\mathrm{out}}}{\partial a_v^{\mathrm{S}}}
      &=\frac{(a_v^{\mathrm{S}}-\bar a^{\mathrm{S}})
              -(a_v^{\mathrm{T}}-\bar a^{\mathrm{T}})}
              {cT^2}+O(T^{-3}).
  \end{aligned}
  \label{eq:ket-integration-temperature-square}
\end{equation}
因此常用$T^2L_{\mathrm{out}}$或$T^2D_{\mathrm{KL}}$，补偿高温区间软目标梯度约按$1/T^2$缩小的效应，使它与$T=1$的真实标签损失在调温时保持大致可比。它并非让所有温度下的梯度完全相同，也不能消除温度改变监督内容的影响。若学生初始分数为零，在例~\ref{ex:ket-integration-soft-labels}~的$T=2$条件下，月季分数的未缩放梯度约为$(1/3-0.5065)/2=-0.0866$，乘$T^2$后约为$-0.3463$，方向仍是提高月季分数。

行为输出也可以是动作分布。Policy Distillation用教师策略或由动作价值转换得到的分布监督学生，并比较硬动作、价值回归与分布匹配。动作编号、可用动作和状态观测必须对应；若在新任务中“动作一”含义改变，同长度分布也不能直接蒸馏。该研究的动作价值分布还可能用低温锐化以突出动作差异，不能把分类中的“提高温度”机械移用到所有策略蒸馏\scite{ket-rusu2015policy}

即使动作接口相同，教师轨迹覆盖的状态也未必覆盖学生执行时会到达的状态；完整策略、探索与价值分析属于第~\ref{part:reinforcement-learning}~部分的内容范围。分类中同样存在教师错误和任务条件丢失：若同一图像在“识别物种”和“判断病害”两种条件下需要不同答案，却不给学生条件，矛盾监督不会因蒸馏消失。较低蒸馏损失只说明更接近选定教师信号，目标准确率仍需真实标签检验。

\subsubsection{基于中间表示的蒸馏}
\label{sec:ket-integration-feature-distillation}

最终类别约束有时不足以有效训练一个较深或较窄的学生。FitNets让教师某个隐藏层输出成为提示，训练学生的一个中间层去预测这个提示。教师宽、学生窄时不能直接逐元素比较，因此引入可学习映射，并明确它属于学生侧训练对象\scite{ket-romero2015fitnets}

设教师第$\ell_{\mathrm{T}}$层输出$\bm{z}^{\mathrm{T}}\in\mathbb{R}^{d_{\mathrm{T}}}$，学生第$\ell_{\mathrm{S}}$层输出$\bm{z}^{\mathrm{S}}\in\mathbb{R}^{d_{\mathrm{S}}}$，映射$r_{\bm{\phi}}:\mathbb{R}^{d_{\mathrm{S}}}\to\mathbb{R}^{d_{\mathrm{T}}}$，则一个提示目标为
\begin{equation}
  L_{\mathrm{rep}}(\bm{x})
   =\frac{1}{2}\left\|
      r_{\bm{\phi}}(\bm{z}^{\mathrm{S}}(\bm{x}))
        -\bm{z}^{\mathrm{T}}(\bm{x})
     \right\|_2^2.
  \label{eq:ket-integration-feature-loss}
\end{equation}
卷积表示还须对应空间位置、通道维度和感受野，可用卷积映射调整，而不能只把张量展平后忽略位置。教师提示的尺度、激活形式和学生映射输出也应可比；若必须插值或聚合位置，需说明对应关系和丢失的信息。

取教学表示$\bm{z}^{\mathrm{T}}=(1,0,2)^\top$、$\bm{z}^{\mathrm{S}}=(0.5,1)^\top$，令
\[
  r_{\bm{\phi}}(\bm{z}^{\mathrm{S}})=\bm{R}\bm{z}^{\mathrm{S}},
  \qquad
  \bm{R}=\begin{bmatrix}2&0\\0&1\\0&2\end{bmatrix}.
\]
同一输入先经过冻结教师得到三维提示，再经过学生得到二维表示，映射输出为$(1,1,2)^\top$，残差为$(0,1,0)^\top$，所以$L_{\mathrm{rep}}=1/2$。映射梯度为$(\bm{R}\bm{z}^{\mathrm{S}}-\bm{z}^{\mathrm{T}})(\bm{z}^{\mathrm{S}})^\top$，学生表示梯度为$\bm{R}^{\top}(\bm{R}\bm{z}^{\mathrm{S}}-\bm{z}^{\mathrm{T}})$；梯度进入映射及学生的前部，教师参数不更新。若暂时只更新映射，学习率取$0.1$，第二行变成$(-0.05,0.9)$，该坐标输出为$0.875$，损失降为$0.3828125$。这一计算也说明，表示维度不同不是把较短向量补零即可解决的问题。

FitNets原始训练先用提示目标训练学生到所选中间层及其映射，再用输出蒸馏目标训练整个学生；也可以设计联合目标，但应说明这是训练安排的变化。只用于提示的映射在部署时通常可以移除，因为最终预测仍沿学生自己的输出路径进行。

精确逐元素匹配要求学生采用与教师相近的内部坐标；允许学习映射，则只要求学生表示足以预测教师提示，约束更灵活。第~\ref{sec:ket-integration-alignment}~节的置换反例已经表明，内部坐标不是语义本身。把过多层、过高维提示都强制匹配，可能限制窄学生适合自身结构的表示。容量过大的映射还可能记住有限训练集上学生表示与教师提示的对应；训练提示损失很低，仍不足以说明学生表示能在未见输入上支持这些信息的提取。因此应在独立验证集上检查去掉映射后的任务表现，而不只看训练中的表示距离。跨结构或跨模态仍沿用同一原则，只是对应和可观测性条件更加严格。

\subsection{承接模型的训练}
\label{sec:ket-integration-student-training}

约束确定后，还需要一个完整训练闭环。设$S_{\mathrm{L}}$为有真实标签的训练集，$U$为蒸馏输入集，二者可以有交集，但都不包含最终测试材料。固定教师与其聚合规则，学生参数为$\bm{\theta}$，提示映射参数为$\bm{\phi}$。一种联合目标为
\begin{equation}
  \begin{aligned}
    L(\bm{\theta},\bm{\phi})={}&
      \frac{\lambda_{\mathrm{L}}}{|S_{\mathrm{L}}|}
      \sum_{(\bm{x},y)\in S_{\mathrm{L}}}
        -\log p_{\bm{\theta}}(y\mid\bm{x};1)\\
      &+\frac{\lambda_{\mathrm{D}}T^2}{|U|}
      \sum_{\bm{x}\in U}
        D_{\mathrm{KL}}\!\left(
          \bm{q}^{(T)}(\bm{x})\Vert
          \bm{p}_{\bm{\theta}}^{(T)}(\bm{x})\right)\\
      &+\frac{\lambda_{\mathrm{R}}}{|U|}
      \sum_{\bm{x}\in U}L_{\mathrm{rep}}(\bm{x};\bm{\theta},\bm{\phi}).
  \end{aligned}
  \label{eq:ket-integration-joint-distillation}
\end{equation}
各权重非负；没有真实标签时删去第一项，不能对空集作除法；不能访问隐藏表示时删去第三项。教师和学生的温度在软目标项中一致，真实标签项按温度一计算。监督权重、温度和学生容量依据允许的验证材料选择，而非依据测试标签选择。学生容量太小时，多个教师之间即使不冲突，也可能没有一个学生参数能同时拟合全部目标。

\begin{algorithm}[从冻结教师到可部署学生]
  \label{alg:ket-integration-distillation-loop}
  \begin{algorithmic}[1]
    \Require 已确定的教师与输出接口，$S_{\mathrm{L}}$、$U$，独立验证集，温度与损失权重，训练预算
    \Ensure 冻结后的学生参数及验证记录
    \State 初始化学生及必要的提示映射；冻结完整教师计算路径
    \For{学习步骤$t=0,1,\ldots$，直到预算耗尽或满足停止条件}
      \State 按预定任务及类别采样比例抽取输入小批次
      \State 查询教师或读取同版本缓存，取得软目标及可用的中间提示
      \State 计算学生的温度一输出、软化输出和必要的映射表示
      \State 按式~\eqref{eq:ket-integration-joint-distillation}计算各可用损失
      \State 只更新学生和提示映射，教师不参与梯度更新
      \State 定期检查验证损失、逐任务保留效果和停止条件，保存合适检查点
    \EndFor
    \State 固定所选学生，移除仅训练需要的教师与映射
    \State 在独立测试材料上检查任务效果、来源保留、组合能力与部署成本
  \end{algorithmic}
\end{algorithm}

算法~\ref{alg:ket-integration-distillation-loop}中的输入采样不应让常见任务淹没少见任务。若缓存教师输出，还须记录教师版本、输入预处理和温度；当数据增强改变输入时，旧缓存未必仍是正确监督。停止依据可以是预算、验证集不再改善或某项保留要求开始恶化，不能只等待训练蒸馏损失接近零。部署通常用温度一输出；若另做概率校准，应把校准参数和所用材料单独记录。

比较学生时，应固定任务、输入范围与评价样本，同时报告原教师、相同结构但只用真实标签训练的学生，以及蒸馏学生。大模型到小模型的转换主要受容量限制，异构网络还受表示匹配和优化条件限制；教师只能通过远程接口访问时，查询预算、返回精度与隐藏层不可见也会改变可行目标。这些条件决定能承接哪些行为，并不因为训练目标写成一个KL项而自动满足。

底座升级是这种跨载体承接的具体情形。Trans-LoRA不把旧底座的LoRA矩阵直接装到新底座上，而是用“旧底座加旧适配器”产生目标输出，再训练新底座上的适配器。其近乎无原数据的迁移仍保留少量任务种子示例，原论文实验使用五个示例，并在原任务训练阶段另训练判别器；迁移时生成候选输入，按判别器筛选，再查询旧模型取得输出。原始完整任务数据不在迁移时重新读取，不意味着生成、筛选、教师查询和学生训练都不需要数据或成本，也不意味着迁移一定无损\scite{ket-wang2024translora}

持续使用蒸馏还会把误差带到下一轮。Progress \& Compress用一个知识库网络保存过去经验，当前任务先由活动分支学习，并通过侧向连接利用冻结知识库的表示；压缩阶段再让活动分支充当教师，把新行为蒸馏回知识库，同时用在线弹性权重巩固约束重要参数。它以固定大小的两部分交替承担当前学习和长期压缩，而不是无限增加任务分支\scite{ket-schwarz2018progress}

压缩时冻结教师，应包括其依赖的旧版本计算或相应缓存，不能一边更新学生，一边无意改变教师目标。保留约束可以减轻历史退化，却不能补回未覆盖的查询区域或无限扩大知识库容量。若每次压缩都只覆盖常见花卉，少见物种的偏差可能逐轮累积；新学生成为下一轮教师后，错误监督还可能被再次复制。应逐任务检查保留效果、通用能力和后续学习速度，并记录教师存储、合成样本、查询和重复训练的费用。第~\ref{sec:ket-update-retention-validation}~节在这一基本过程上补充持续更新的验证；模型压缩算子与服务部署优化分别属于模型卷及第~\ref{part:systems}~部分。

\section{本章小结}
\label{sec:ket-integration-summary}

分别训练的来源能够怎样共同发挥作用，取决于能力进入结果的接口。预测集成要求答案可比，收益取决于错误互补和组合规则；模块组合要求中间语义、尺度与状态能够衔接，局部组件正确不足以保证整条链正确。固定与条件描述的是任务或输入变化时组合是否改变，与组合器是否经过训练是两件事。

参数融合把多来源计算压到一份参数中，代价是必须处理坐标对应、非线性和有限容量带来的冲突。隐藏单元置换说明函数相同也可能不能直接逐坐标平均，LoRAHub的交叉项说明因子组合也不是完整增量平均。蒸馏放宽了参数和架构对应的要求，却把困难转移到教师可靠性、输入覆盖、监督可比性和学生训练上；较低匹配损失不代表更多真实知识或更准确的目标行为。

回到花卉系统，能够承担多来源推理开销时可以保留组合；参数关系足够明确时可以尝试融合；希望更换结构或压缩部署时可以利用教师训练学生，也可以把组合或融合结果继续蒸馏。无论选择哪条路径，都应分别检验来源能力保留与未见能力组合，并核算一次整合和长期维护的成本。第~\ref{chap:ket-transfer-adaptation}~章进一步研究如何为新目标形成可用行为，第~\ref{chap:ket-sharing}~章研究多个任务共同学习时怎样形成可共享知识，统一的证据与评价协议见第~\ref{chap:ket-evaluation}~章。
